\section{Weight access against function access}\label{sec:function}

Corollary~\ref{cor:ceiling} says a \emph{complete} $G$-invariant of the weights carries exactly the
information of the realized function, so weight access can only win on compute. Function access
evaluates $f_\theta$ at $K$ \emph{learned} probe coordinates and classifies the outputs; nothing
reads a weight, and our test suite certifies the reader is exactly $G$-invariant. Learning the
probes is the strong form of the baseline, which is the conservative direction here, and FLOPs are
analytic per INR at inference.

At $K = 64$, function access reaches $95.34\%$ for $1.59$\,MFLOP per INR against $\calign$'s
$64.41\%$ for $5.45$\,MFLOP, and at $K = 256$ it reaches $98.23\%$, above the real-pixel MLP's
$97.97\%$, so learned query points are a better input to the same decoder than the pixels are. W12
did not exist when this was registered and the claim is quantified over all weight pipelines, so it
has to be re-priced: W12 reaches $87.64\%$ at $163$\,MFLOP, which is $7.7$ accuracy points worse at
$103\times$ the compute. The conclusion survives, but its shape changes, from a large accuracy
deficit at comparable cost to a small one at two orders of magnitude more compute. Amortizing does
not rescue it. Weight access costs $1.70 + 3.74\,T$\,MFLOP over $T$ downstream tasks against
function-query's $1.59\,T$, and the lines never cross, because the weight reader's \emph{per-task}
cost on a $1185$-dimensional input already exceeds function-query's \emph{entire} per-task cost on a
$64$-dimensional one. The mechanism is general: function access is cheaper whenever $Kc \ll P$, a
condition on how many probes the task needs rather than on how big the INR is.

One asymmetry is worth stating on its own. Between \texttt{P-random} and \texttt{P-shared-det},
function-query accuracy moves $5.4$ points, a fit-quality effect, while weight access moves $80.4$.
The entire object this paper decomposes is invisible from the function side. The scope is that the
result holds where the function is cheap to query relative to the parameter count, for targets that
are functions of the realized signal, at width $32$ and depth $2$; what it settles is the case that
has been made for weight-space learning on INR benchmarks of this kind.
